\section{Introducción}

El consumismo ha caracterizado a la sociedad desde finales de la Segunda
Guerra Mundial. En ese contexto, las necesidades básicas de la población ya
estaban cubiertas; sin embargo, se continuó priorizando el crecimiento
económico, para lo cual se recurrió a campañas de marketing masivas que
manipulaban a los consumidores y generaban de forma constante nuevas
necesidades \parencite[p.~160]{gorostiza2025}. Hoy en día esta dinámica
persiste: las campañas siguen haciendo creer al consumidor que necesita
determinados productos, y es ahí donde interviene la obsolescencia programada,
que perpetúa el mismo ciclo de consumo. Esta investigación analiza cómo ha
cambiado el comportamiento de la sociedad a raíz de los avances tecnológicos y
de crisis globales como el COVID-19, las cuales evidencian la vulnerabilidad del
ser humano y la facilidad con que puede ser manipulado al tomar decisiones.

Con base en lo anterior, la presente investigación busca informar y respaldar el
propósito de la aplicación ``TreeBar'', que pretende transformar tanto los
hábitos de consumo como la forma de pensar del consumidor. TreeBar centraliza la
información de los productos que el usuario desea comprar y, además, hace
visibles los procesos de producción y administrativos de las empresas: cómo y
dónde fabrican sus productos y qué impacto tienen sus decisiones en el ámbito
sociopolítico. De este modo, al conocer las decisiones que toma cada empresa en
el mercado, el usuario podrá saber realmente qué tipo de negocio está apoyando
con su compra.

Asimismo, la aplicación ofrecerá opiniones auténticas de compradores reales.
Cada producto contará con su propio apartado de reseñas, dedicado
exclusivamente a experiencias de uso real y cotidiano. Además, quienes deseen
adquirir un producto podrán formular preguntas sobre aspectos que solo se
aprecian con el uso normal, las cuales serán respondidas por quienes ya hayan
escrito una reseña.

Otra funcionalidad práctica será la recomendación de lugares donde reparar
productos, como smartphones o laptops, de modo que el usuario que aprecie su
dispositivo pueda repararlo en establecimientos respaldados por los comentarios
de personas que ya han acudido a ellos. Todo ello refuerza lo expuesto a lo
largo de la investigación y da sentido a TreeBar.